\chapter{Expected Outcomes}
\label{ch:outcomes}

Each research question can come out in more than one way. Table~\ref{tab:out-map}
fixes in advance how each outcome will be interpreted, so that negative results
can be read as clearly as positive ones. Each module is measured on its own part of the cost,
so a module that saves little does not undermine the others, and the cost
breakdown of RQ1 is a result in any case.

\begin{table}[htbp]
\centering
\caption{Possible outcomes and how they will be read.}
\label{tab:out-map}
\footnotesize
\begin{tabularx}{\textwidth}{@{}>{\raggedright\arraybackslash}p{0.07\textwidth} >{\raggedright\arraybackslash}p{0.38\textwidth} >{\raggedright\arraybackslash}X@{}}
\toprule
RQ & If \dots & then \dots\\
\midrule
RQ1 & validation is a large share of the cost & the Repeat Guard (M1) and the Test Prioritiser (M2) matter; the headroom says how much they can save.\\
 & generation dominates & the repeat notice and the Budget Guard are the main levers.\\
 & little cost falls after the model is stuck & the Budget Guard is not needed in this setting.\\
\midrule
RQ2 & \km\ is cheaper than \frtp\ and Coverage & learning failure probability and cost during the session is effective.\\
 & \km\ only matches Coverage & its remaining advantage is needing no instrumentation; Coverage is preferred where coverage is available.\\
 & \km\ saves on mutants (E1) but not live (E2) & the feedback it reports changes the model's patches; this is reported as a finding.\\
 & an ordering lowers the repair rate & the failure it shows helps the model less; the saving is weighed against the loss.\\
 & with a trigger step, ordering the rest saves little & the practical value of the Test Prioritiser (M2) is small; the cost breakdown and the mechanism of E1 still stand.\\
 & the Budget Guard saves much at small loss & stopping unproductive sessions reduces cost at little loss.\\
\midrule
RQ3 & the layer costs less with the same repair rate & the layer meets the need; the breakdown shows which module saved what.\\
 & the layer costs less but repairs fewer bugs & the trade-off is reported, and E2 shows which module caused it.\\
 & no saving end to end & the savings are diluted by costs the layer does not touch; RQ1 shows which.\\
\midrule
RQ4 & one part of \km\ explains most of its gain & that part is the recommendation; the rest is not needed.\\
 & the gain disappears with a wrong location & \km\ depends on knowing where the patch is.\\
 & ordering by the posterior mean does as well as sampling & exploration does not pay within a session; the simpler greedy order is recommended.\\
\bottomrule
\end{tabularx}
\end{table}

The layer is expected to help most where test times differ widely, sessions
have many wrong patches, and the model repeats itself often
(Section~\ref{sec:sys-regimes}).
